% ============================================================
%  WTB Workload Engine — Technical Architecture Document
%  Compile with: pdflatex architecture.tex (twice for TOC)
% ============================================================
\documentclass[12pt,a4paper]{report}

% ── Packages ──────────────────────────────────────────────
\usepackage[T1]{fontenc}
\usepackage[utf8]{inputenc}
\usepackage[english]{babel}
\usepackage{lmodern}
\usepackage{geometry}
\geometry{margin=2.5cm, top=3cm, bottom=3cm}

\usepackage{amsmath, amssymb, amsthm}
\usepackage{mathtools}
\usepackage{booktabs}
\usepackage{tabularx}
\usepackage{longtable}
\usepackage{array}
\usepackage{xcolor}
\usepackage{listings}
\usepackage{graphicx}
\usepackage{hyperref}
\usepackage{enumitem}
\usepackage{fancyhdr}
\usepackage{titlesec}
\usepackage{tcolorbox}
\usepackage{microtype}
\usepackage{float}

% ── Colours ───────────────────────────────────────────────
\definecolor{primary}{RGB}{30, 80, 160}
\definecolor{codebg}{RGB}{245, 245, 248}
\definecolor{keywordcolor}{RGB}{0, 100, 180}
\definecolor{commentcolor}{RGB}{100, 120, 100}
\definecolor{stringcolor}{RGB}{160, 50, 50}
\definecolor{highred}{RGB}{200, 40, 40}
\definecolor{medorange}{RGB}{210, 120, 0}
\definecolor{lowgreen}{RGB}{40, 160, 80}

% ── Hyperref setup ────────────────────────────────────────
\hypersetup{
  colorlinks=true,
  linkcolor=primary,
  urlcolor=primary,
  citecolor=primary,
  pdftitle={WTB Workload Engine -- Technical Architecture},
  pdfauthor={Engineering Team},
}

% ── Headers ───────────────────────────────────────────────
\pagestyle{fancy}
\fancyhf{}
\fancyhead[L]{\small\textcolor{gray}{WTB Workload Engine}}
\fancyhead[R]{\small\textcolor{gray}{\leftmark}}
\fancyfoot[C]{\thepage}
\renewcommand{\headrulewidth}{0.4pt}

% ── Code listings ─────────────────────────────────────────
\lstset{
  backgroundcolor=\color{codebg},
  basicstyle=\ttfamily\small,
  keywordstyle=\color{keywordcolor}\bfseries,
  commentstyle=\color{commentcolor}\itshape,
  stringstyle=\color{stringcolor},
  breaklines=true,
  frame=single,
  framerule=0.4pt,
  rulecolor=\color{gray!50},
  showstringspaces=false,
  tabsize=4,
  captionpos=b,
}

% ── Theorem-like boxes ────────────────────────────────────
\newtheorem{definition}{Definition}[chapter]
\newtheorem{property}{Property}[chapter]

\tcbset{
  defbox/.style={
    colback=blue!5, colframe=primary,
    fonttitle=\bfseries, title={Definition},
    sharp corners, left=4pt, right=4pt,
  },
  notebox/.style={
    colback=yellow!10, colframe=orange!70!black,
    fonttitle=\bfseries,
    sharp corners, left=4pt, right=4pt,
  },
}

% ── Math operators ────────────────────────────────────────
\DeclareMathOperator{\vh}{VH}
\newcommand{\floor}[1]{\left\lfloor #1 \right\rfloor}
\newcommand{\ceil}[1]{\left\lceil #1 \right\rceil}

% ── Section formatting ────────────────────────────────────
\titleformat{\chapter}[display]
  {\normalfont\huge\bfseries\color{primary}}
  {\chaptertitlename\ \thechapter}{16pt}{\Huge}

% ============================================================
\begin{document}

% ── Title page ────────────────────────────────────────────
\begin{titlepage}
  \centering
  \vspace*{2cm}
  {\huge\bfseries\color{primary} WTB Workload Engine}\\[0.6cm]
  {\Large Technical Architecture \& Algorithm Reference}\\[0.3cm]
  \rule{\linewidth}{0.5pt}\\[1cm]
  {\large Engineering Document — Version 1.0}\\[0.5cm]
  {\large May 2026}\\[3cm]

  \begin{tcolorbox}[notebox, title={Scope of this document}]
    This document describes the complete architecture, data model, and
    algorithms of the WTB Workload Engine backend. It is written for
    engineers who need to understand, maintain, or extend the system.
    Frontend implementation details are out of scope; the analytical
    engine is the central subject.
  \end{tcolorbox}

  \vfill
  {\small\textcolor{gray}{Internal document --- WTB Qualification Team}}
\end{titlepage}

% ── Table of contents ─────────────────────────────────────
\tableofcontents
\clearpage

% ============================================================
\chapter{Context and Objectives}
% ============================================================

\section{Business Problem}

The WTB (Work-to-be-done) qualification process covers a portfolio of
\textbf{589 topics}---each representing one supplier product or component
to be qualified---spread across \textbf{11 active Certification Agents}
(CAs). Every topic goes through five sequential qualification phases.
Each phase carries its own time budget and labour cost.

Without tooling, a manager has no visibility into:
\begin{itemize}
  \item which CAs are currently overloaded or underutilised;
  \item whether a workload peak is approaching in the coming months;
  \item which specific topics should be reassigned to rebalance the team.
\end{itemize}

The WTB Workload Engine solves this by computing, from a structured
Excel workbook, a full 78-week workload forecast per CA and generating
concrete rebalancing suggestions and actionable recommendations.

\section{System Overview}

The system consists of two independent processes:

\begin{enumerate}
  \item \textbf{Backend (Python, port 8000)} --- a FastAPI application
        that hosts the analytical engine. All business logic, calculations,
        and data transformations live exclusively here.
  \item \textbf{Frontend (TypeScript/React, port 5173)} --- a single-page
        application that calls the backend REST API and renders results.
        The frontend contains no business logic; it is a pure view layer.
\end{enumerate}

\begin{tcolorbox}[notebox, title={Design principle}]
  Every algorithm described in this document lives in
  \texttt{backend/domain/engine.py}, which is a module of pure functions
  with no I/O and no framework dependencies. This makes the engine fully
  testable in isolation and independent of any web framework.
\end{tcolorbox}

% ============================================================
\chapter{Input Data Model}
% ============================================================

\section{The Excel Workbook}

All data enters the system through a single \texttt{.xlsx} workbook
uploaded by the user. The workbook must contain exactly four sheets.
Any structural deviation causes a descriptive validation error before
computation begins.

\subsection{Sheet: \texttt{config}}

A two-column key--value table that parameterises the entire engine.
No hard-coded constants appear anywhere in the algorithms; all tunable
values come from this sheet.

\begin{table}[H]
\centering
\caption{Configuration parameters and their roles}
\begin{tabularx}{\linewidth}{@{} l l X @{}}
\toprule
\textbf{Key} & \textbf{Example} & \textbf{Role in the engine} \\
\midrule
\texttt{today} & 2026-04-24 &
  The reference date for all temporal computations. Using a
  configurable date (rather than the system clock) makes the engine
  fully reproducible: the same file always yields the same results. \\[4pt]
\texttt{weekly\_capacity\_hrs} & 40 &
  The maximum hours a CA can deliver in one week. Used as the
  denominator in every workload-fraction computation. \\[4pt]
\texttt{catchup\_window\_weeks} & 4 &
  Number of weeks into which overdue work is compressed. See
  Section~\ref{sec:effective-window}. \\[4pt]
\texttt{stale\_threshold\_weeks} & 12 &
  A phase whose scheduled end is more than this many weeks before
  \texttt{today} and that has not been closed is flagged as stale
  and excluded from forecast. \\[4pt]
\texttt{forecast\_horizon\_weeks} & 78 &
  How far into the future the engine projects workload ($\approx$18
  months). \\
\bottomrule
\end{tabularx}
\end{table}

\subsection{Sheet: \texttt{phase\_params}}

One row per phase. Defines the intrinsic cost of each qualification phase.

\begin{table}[H]
\centering
\caption{Phase parameter columns}
\begin{tabular}{@{} l l l @{}}
\toprule
\textbf{Column} & \textbf{Type} & \textbf{Meaning} \\
\midrule
\texttt{phase}         & string  & Phase identifier (index key) \\
\texttt{duration\_hrs} & float   & Total labour hours the phase requires \\
\texttt{delai\_weeks}  & float   & Standard calendar duration in weeks \\
\bottomrule
\end{tabular}
\end{table}

The five phases in order are:
\texttt{advanced}, \texttt{digital}, \texttt{static}, \texttt{dynamic},
\texttt{delivery}.

\subsection{Sheet: \texttt{persons}}

One row per CA (both active and inactive). Key columns:

\begin{table}[H]
\centering
\begin{tabular}{@{} l l @{}}
\toprule
\textbf{Column} & \textbf{Meaning} \\
\midrule
\texttt{ca\_sheet}   & Unique identifier used as foreign key in \texttt{topics} \\
\texttt{ca\_display} & Human-readable name shown in the UI \\
\texttt{active}      & 1 = included in all computations; 0 = ignored \\
\texttt{n\_topics}   & Pre-computed count of topics assigned to this CA \\
\bottomrule
\end{tabular}
\end{table}

\subsection{Sheet: \texttt{topics}}

The main fact table. One row per topic. Columns include:

\begin{itemize}
  \item \texttt{id} --- unique integer topic identifier.
  \item \texttt{ca\_sheet} --- foreign key into \texttt{persons}.
  \item \texttt{project}, \texttt{topic}, \texttt{supplier} --- descriptive
        labels.
  \item \texttt{global\_status} --- the overall qualification lifecycle
        status of the topic.
  \item For each of the five phases $p \in \{\texttt{advanced},\ldots\}$:
    \begin{itemize}
      \item \texttt{$p$\_start}, \texttt{$p$\_end} --- scheduled date range
            (may be blank if the phase is not yet planned).
      \item \texttt{$p$\_status} --- current execution status for that
            phase.
    \end{itemize}
\end{itemize}

This gives $3 \times 5 = 15$ phase-specific columns per row, plus
general columns, yielding approximately 20 columns per topic.

\subsection{Parsing and Validation}

The function \texttt{parse\_workbook(bytes) $\to$ DataBundle} in
\texttt{backend/domain/loader.py} performs:

\begin{enumerate}
  \item Load all four sheets with \texttt{openpyxl}.
  \item Assert all required sheets and columns are present.
  \item Parse the 10 date columns in \texttt{topics} directly to
        \texttt{pd.Timestamp} objects.
  \item Index \texttt{phase\_params} by phase name for $O(1)$ lookup.
  \item Coerce \texttt{config} values to native Python scalars or
        \texttt{pd.Timestamp}.
\end{enumerate}

The result is a \texttt{DataBundle} dataclass with four fields:
\texttt{topics}, \texttt{phase\_params}, \texttt{persons}, \texttt{config}.
This is the single immutable input to every downstream computation.

% ============================================================
\chapter{The Analytical Engine}
% ============================================================

\section{Introduction}

The engine transforms a \texttt{DataBundle} into workload fractions
through a pipeline of three major stages:

\begin{enumerate}
  \item \textbf{Phase instance construction} --- flatten the topic table
        into one row per (CA, topic, phase) triple, applying temporal
        adjustments and classification flags.
  \item \textbf{Workload forecast} --- aggregate phase instances into
        a $(W \times C)$ matrix of weekly workload fractions over the
        forecast horizon, where $W$ is the number of weeks and $C$ is
        the number of active CAs.
  \item \textbf{Higher-level outputs} --- derive KPIs, smoothing
        suggestions, and prioritised recommendations from the forecast
        matrix.
\end{enumerate}

All three stages are pure deterministic functions. Given the same
\texttt{DataBundle} they always produce the same output.

% ────────────────────────────────────────────────────────────
\section{Stage 1 --- Phase Instance Construction}
\label{sec:phase-instances}
% ────────────────────────────────────────────────────────────

\subsection{Overview}

\texttt{build\_phase\_instances(bundle) $\to$ DataFrame}

The topics sheet stores dates in a wide format (one topic per row,
phases in columns). The first step \textit{normalises} this into a long
format: one row per active-CA $\times$ topic $\times$ phase triple.
Each row represents an atomic unit of work and carries all the
information needed to place it on the forecast timeline.

\subsection{Phase Intensity (VH)}
\label{sec:vh}

Before computing any windows, the engine derives the \textbf{weekly
workload intensity} of each phase, called $\vh$.

\begin{tcolorbox}[defbox, title={Definition: Workload Intensity $\vh$}]
\[
  \vh_p \;=\; \frac{\texttt{duration\_hrs}_p}{\texttt{delai\_weeks}_p}
\]
where $p$ denotes one of the five qualification phases.
\end{tcolorbox}

$\vh_p$ expresses how many hours per week a CA must spend on phase $p$
while it is active. It is constant per phase type, independent of
any specific topic or date.

\textbf{Why this formulation?}
The system models each phase as a uniform workload rectangle: $\vh_p$
hours per week, sustained over the effective duration of the phase.
This is an engineering approximation that avoids speculative
time-phasing of individual tasks. It is well-suited to planning
horizons of several months, where daily granularity would be spurious.

\begin{table}[H]
\centering
\caption{Example phase parameters and resulting VH values (illustrative)}
\begin{tabular}{@{} l r r r @{}}
\toprule
\textbf{Phase} & \textbf{duration\_hrs} & \textbf{delai\_weeks} & $\vh$ (h/week) \\
\midrule
advanced  & 24 & 4 & 6.0 \\
digital   & 16 & 2 & 8.0 \\
static    & 40 & 6 & 6.7 \\
dynamic   & 32 & 4 & 8.0 \\
delivery  &  8 & 1 & 8.0 \\
\bottomrule
\end{tabular}
\end{table}

\subsection{The Effective Window}
\label{sec:effective-window}

A key challenge is that many phases in a real portfolio are overdue:
their scheduled end date has already passed, yet the work has not
been completed or formally closed. Simply ignoring these phases
would undercount the actual current workload. But placing them at
their original (past) dates would also be wrong, since the work will
be done now, not in the past.

The engine resolves this with the \textbf{effective window} concept.
Given the scheduled start $t_s$ and end $t_e$ of a phase instance,
and a reference date $t_0$ (``today''), the effective window
$[t_s^*, t_e^*]$ is computed as follows.

\begin{tcolorbox}[defbox, title={Definition: Effective Window}]
Let $\Delta_c$ denote the catchup window in weeks
(\texttt{catchup\_window\_weeks}).

\medskip
\textbf{Effective start:}
\[
  t_s^* \;=\; \max(t_0,\; t_s)
\]
If $t_s$ is absent, $t_s^* = t_0$. This ensures that future
phases begin no earlier than scheduled, while past or present
phases begin immediately from today.

\medskip
\textbf{Catchup end:}
\[
  t_c \;=\; t_s^* + \Delta_c \text{ weeks}
\]

\textbf{Effective end:}
\[
  t_e^* \;=\; \max(t_c,\; t_e)
\]
If $t_e$ is absent, $t_e^* = t_c$.
\end{tcolorbox}

\textbf{Interpretation of these rules:}

\begin{enumerate}
  \item If a phase is \textit{future} ($t_s \geq t_0$): the effective
        window equals the scheduled window (possibly extended to $t_c$
        if the scheduled duration is shorter than $\Delta_c$).

  \item If a phase is \textit{overdue} ($t_s < t_0$ or $t_e < t_0$):
        the effective start collapses to today, and the work is spread
        over the catchup window $\Delta_c = 4$ weeks starting from today.
        This models the operational reality that overdue work
        becomes urgent and must be completed in the near term.

  \item The $\max(t_c, t_e)$ in the effective end ensures that a phase
        that is both overdue \textit{and} has a far future end date is
        not artificially compressed.
\end{enumerate}

\begin{figure}[H]
\centering
\begin{tcolorbox}[colback=gray!8, colframe=gray!50, sharp corners]
\small
\begin{verbatim}
  Case 1 (future phase):
  ─────────────────────────────────────────────────────
  t0          ts          te
  |           [===PHASE===]
              t_s*        t_e* (= te if te > ts* + Δ_c)

  Case 2 (overdue phase, te < t0):
  ─────────────────────────────────────────────────────
       te   t0
  [====]    |
            [==CATCHUP==]
            t_s*         t_e* = t_s* + Δ_c

  Case 3 (started, not yet ended, t_s < t0 < t_e):
  ─────────────────────────────────────────────────────
       ts   t0            te
  [====|====|=============]
            t_s*  max(t_s* + Δ_c, te)
                          t_e*
\end{verbatim}
\end{tcolorbox}
\caption{Effective window construction for the three cases}
\end{figure}

\subsection{Stale Phase Detection}
\label{sec:stale}

A phase that ended far in the past and was never formally closed is
likely an administrative artefact rather than real pending work.
Including it would inflate every CA's workload permanently.

\begin{tcolorbox}[defbox, title={Definition: Stale Phase}]
A phase instance is \textbf{stale} if and only if:
\[
  t_e < t_0 - \Delta_s \;\text{ weeks}
  \quad \text{and} \quad
  \texttt{status} \notin \{\texttt{VAL},\, \texttt{CANCELED}\}
\]
where $\Delta_s$ = \texttt{stale\_threshold\_weeks} (currently 12 weeks).
\end{tcolorbox}

Stale phases are retained in the dataset with an \texttt{is\_stale = True}
flag but are \textit{excluded from all workload computations}.
They remain visible in the CA detail view so a manager can investigate
them as administrative tasks, without them polluting the forecast.

\begin{tcolorbox}[notebox, title={Note on VAL and CANCELED}]
Phases with status \texttt{VAL} (validated, work completed) or
\texttt{CANCELED} are excluded unconditionally---regardless of their
dates---from every workload computation. These statuses represent
definitively closed work.
\end{tcolorbox}

\subsection{Active-Now Classification}

After building the effective window for every phase instance, the engine
classifies each row with a boolean \texttt{is\_active\_now} flag.

\begin{tcolorbox}[defbox, title={Definition: Active-Now Phase}]
Let $[w_0, w_1)$ be the current week interval:
\[
  w_0 = \text{Monday of the week containing } t_0,
  \qquad
  w_1 = w_0 + 7 \text{ days}
\]
A phase instance is \textbf{active now} if all four conditions hold:
\begin{align*}
  &(1)\quad \texttt{status} \notin \{\texttt{VAL}, \texttt{CANCELED}\} \\
  &(2)\quad \texttt{is\_stale} = \texttt{False} \\
  &(3)\quad t_s^* < w_1 \\
  &(4)\quad t_e^* \geq w_0
\end{align*}
\end{tcolorbox}

Conditions (3) and (4) are the standard interval-overlap test:
the effective window $[t_s^*, t_e^*]$ overlaps with $[w_0, w_1)$
if and only if $t_s^* < w_1$ and $t_e^* \geq w_0$.

\texttt{is\_active\_now} drives the KPI ``number of active topics per CA''
and is the basis for identifying which topics to consider for transfer
in the smoothing algorithm.

% ────────────────────────────────────────────────────────────
\section{Stage 2 --- Workload Forecast}
\label{sec:forecast}
% ────────────────────────────────────────────────────────────

\subsection{Weekly Workload Fraction}

\texttt{forecast\_weekly(bundle, phase\_instances) $\to$ DataFrame}

The forecast is a matrix $\mathbf{W} \in \mathbb{R}^{W \times C}$ where:
\begin{itemize}
  \item $W = H + 1$ rows, one per week ($H$ = \texttt{forecast\_horizon\_weeks}).
        Week indices are the Mondays $w_0, w_1, \ldots, w_H$.
  \item $C$ columns, one per active CA, sorted lexicographically.
  \item Each element $\mathbf{W}_{i,c}$ is a \textbf{workload fraction},
        where 1.0 represents exactly full capacity for CA $c$ during week $i$.
\end{itemize}

\begin{tcolorbox}[defbox, title={Definition: Weekly Workload Fraction}]
Let $\mathcal{P}(c, i)$ be the set of non-stale, non-excluded phase
instances assigned to CA $c$ whose effective window overlaps with week
$[w_i,\, w_i + 7\text{ days})$:
\[
  \mathcal{P}(c, i) = \bigl\{
    p \in \Pi :
    \texttt{ca\_id}(p) = c,\;
    t_s^*(p) < w_i + 7\text{d},\;
    t_e^*(p) \geq w_i
  \bigr\}
\]
where $\Pi$ excludes all stale or definitively closed instances.

The workload fraction is:
\[
  \boxed{
    \mathbf{W}_{i,c} \;=\;
    \frac{1}{\texttt{capacity}} \sum_{p \in \mathcal{P}(c,\,i)} \vh_p
  }
\]
\end{tcolorbox}

\textbf{What this formula computes.}
For each week $i$ and each CA $c$, the engine sums the hourly intensity
$\vh_p$ of every phase that is simultaneously assigned to CA $c$ and
overlapping that week's calendar interval. Dividing by the weekly capacity
(40 h) converts the raw sum of hours per week into a dimensionless
fraction, where 1.0 means the CA is working exactly at full capacity.

\textbf{Values above 1.0} are structurally possible and expected: a CA
with three overlapping phases each contributing 15 h/week would produce
$\mathbf{W}_{i,c} = 45/40 = 1.125$, indicating an overload of 12.5\%.

\subsection{Vectorised Implementation}

The naive double-loop over weeks and phase instances would be $O(WN)$
with two passes over the data. The engine uses NumPy broadcasting to
compute the full $(W \times N)$ overlap boolean matrix in a single
vectorised operation:

\begin{lstlisting}[language=Python, caption={Vectorised overlap computation}]
# eff_starts, eff_ends : shape (N,)  -- one per phase instance
# ws_np, we_np         : shape (W,)  -- one per week boundary

overlap = (
    (eff_starts[np.newaxis, :] < we_np[:, np.newaxis])   # (W, N)
  & (eff_ends  [np.newaxis, :] >= ws_np[:, np.newaxis])
)
# overlap[i, j] = True iff phase j overlaps week i
\end{lstlisting}

Broadcasting inserts a new axis so NumPy compares every week against
every phase instance in one operation. For 79 weeks and $\sim$2000 active
phase instances this produces a $79 \times 2000$ boolean matrix
in a single pass. Summation over CA is then a grouped aggregation.

\subsection{Monthly Aggregation}

\texttt{forecast\_monthly(weekly\_df) $\to$ DataFrame}

The monthly view averages the weekly fractions over each calendar month:

\[
  \mathbf{W}^{\text{month}}_{m,c}
  \;=\;
  \frac{1}{|\mathcal{W}_m|}
  \sum_{i \in \mathcal{W}_m} \mathbf{W}_{i,c}
\]

where $\mathcal{W}_m$ is the set of week indices whose Monday falls
in calendar month $m$. This is implemented as a single pandas
\texttt{resample("MS").mean()} call, which groups weeks by their
Month-Start timestamp and averages across the group.

The monthly view smooths short-term noise (e.g.\ a single heavy week)
and gives management a strategic, high-level visibility of workload
trends over the 18-month horizon.

\subsection{Week Series Generation}

All week indices are anchored to the Monday of the week containing
\texttt{today}, ensuring that row 0 of $\mathbf{W}$ always corresponds
to the \textit{current} week, regardless of which day the file is loaded.

\[
  w_0 = t_0 - (\text{weekday}(t_0) \bmod 7) \text{ days}
  \qquad
  w_i = w_0 + 7i \text{ days},\quad i = 0, 1, \ldots, H
\]

% ────────────────────────────────────────────────────────────
\section{Stage 3A --- Key Performance Indicators}
\label{sec:kpis}
% ────────────────────────────────────────────────────────────

From the forecast matrix and phase instances, the engine derives a set
of KPIs for the overview dashboard.

\begin{table}[H]
\centering
\caption{Computed KPIs}
\begin{tabularx}{\linewidth}{@{} l X @{}}
\toprule
\textbf{KPI} & \textbf{Formula / derivation} \\
\midrule
\texttt{total\_active\_topics} &
  $|\{topic\_id : \exists p, \texttt{is\_active\_now}(p) = \texttt{True}\}|$
  --- count of distinct topics with at least one phase active this week. \\[6pt]
\texttt{total\_active\_cas} &
  $C$ --- number of columns in $\mathbf{W}$ (always equals the number
  of CAs with \texttt{active}=1 in \texttt{persons}). \\[6pt]
\texttt{avg\_workload} &
  $\bar{w}_0 = \frac{1}{C} \sum_c \mathbf{W}_{0,c}$
  --- mean workload fraction across all CAs for the current week. \\[6pt]
\texttt{max\_workload} &
  $\max_c \mathbf{W}_{0,c}$ --- workload of the most loaded CA
  this week. \\[6pt]
\texttt{max\_workload\_ca} &
  $\arg\max_c \mathbf{W}_{0,c}$ --- identifier of that CA. \\[6pt]
\texttt{peak\_workload} (per CA) &
  $\max_i \mathbf{W}_{i,c}$ --- highest workload fraction for CA $c$
  over the entire forecast horizon. \\[6pt]
\texttt{peak\_week} (per CA) &
  $\arg\max_i \mathbf{W}_{i,c}$ --- the Monday date of that peak. \\
\bottomrule
\end{tabularx}
\end{table}

\textbf{Active topics per CA.}
A topic counts as active for a CA when at least one of its phases is
currently active-now for that CA. This is a distinct-count of
\texttt{topic\_id} in the subset of rows where \texttt{is\_active\_now =
True} and \texttt{ca\_id = c}.

% ────────────────────────────────────────────────────────────
\section{Stage 3B --- Workload Smoothing}
\label{sec:smoothing}
% ────────────────────────────────────────────────────────────

\subsection{Problem Statement}

When one or more CAs are overloaded this week
($\mathbf{W}_{0,c} > 1.0$), the engine proposes a concrete transfer:
move one specific topic from the overloaded CA to the CA with the
most available capacity that can absorb it.

The algorithm operates on $\mathbf{w}_0 = \mathbf{W}_{0,\cdot}$
(the current-week row of the forecast matrix), treated as a mutable
working copy to simulate the effect of sequential transfers.

\subsection{Algorithm}

\begin{tcolorbox}[defbox, title={Smoothing Algorithm}]
\textbf{Input:} Working vector $\mathbf{w} \in \mathbb{R}^C$
(current-week workload fractions), phase instances, capacity $K$.

\medskip
\textbf{Step 1.} Identify overloaded CAs:
\[
  \mathcal{O} = \{ c : \mathbf{w}_c > 1.0 \}
  \quad \text{sorted descending by } \mathbf{w}_c
\]

\medskip
\textbf{Step 2.} For each $c \in \mathcal{O}$ (if $\mathbf{w}_c \leq 1.0$
after a prior step, skip):

\quad \textbf{2a.} Find all topics with at least one active-now phase
assigned to $c$. For each such topic $\tau$, compute its load
contribution:
\[
  \delta_\tau \;=\;
  \frac{1}{K} \sum_{p \in \Pi(c,\tau)} \vh_p
\]
where $\Pi(c, \tau)$ is the set of active-now phases for CA $c$ and
topic $\tau$. Select the topic with the highest contribution:
\[
  \tau^* = \arg\max_\tau \delta_\tau
\]

\quad \textbf{2b.} Find the best receiver among all other active CAs:
\[
  \mathcal{R}^* = \{c' \neq c : \mathbf{w}_{c'} + \delta_{\tau^*} \leq 1.0\}
\]
If $\mathcal{R}^* = \emptyset$, relax the constraint:
\[
  \mathcal{R}^* = \{c' \neq c : \mathbf{w}_{c'} + \delta_{\tau^*} \leq 1.2\}
\]
If still empty, skip this CA (no feasible transfer).

\quad \textbf{2c.} Select receiver:
\[
  c^* = \arg\min_{c' \in \mathcal{R}^*} \mathbf{w}_{c'}
\]

\quad \textbf{2d.} Record suggestion and update the working vector:
\[
  \mathbf{w}_c \;\leftarrow\; \mathbf{w}_c - \delta_{\tau^*}
  \qquad
  \mathbf{w}_{c^*} \;\leftarrow\; \mathbf{w}_{c^*} + \delta_{\tau^*}
\]
\end{tcolorbox}

\subsection{Output}

Each suggestion contains:
\begin{itemize}
  \item Source and destination CA identifiers and display names.
  \item The topic ID, name, and project.
  \item An \textbf{impact object} showing the before/after workload
        fractions for both the source and the destination CA:
        $(\mathbf{w}_c^{\text{before}},\; \mathbf{w}_c^{\text{after}},\;
           \mathbf{w}_{c^*}^{\text{before}},\; \mathbf{w}_{c^*}^{\text{after}})$.
\end{itemize}

\subsection{Design Choices}

\textbf{One suggestion per overloaded CA.}
The algorithm proposes one transfer per overloaded CA per run. This
intentional constraint prevents the output from overwhelming the manager
with a complex multi-step rebalancing plan. One change at a time is
easier to evaluate and implement.

\textbf{Greedy selection of the heaviest topic.}
Transferring the highest-load topic maximises the immediate load
reduction per transfer. A branch-and-bound approach would find the
globally optimal transfer sequence, but given the small team size
($C \leq 14$) and the qualitative nature of the recommendations,
a greedy approach is entirely adequate.

\textbf{1.2 relaxed bound.}
If no receiver can stay under 1.0 after the transfer, the engine
relaxes to 1.2 (20\% overload). This avoids producing no suggestions
in situations where every CA is heavily loaded---in such cases, even
a partial redistribution is useful to present to the manager.

\textbf{Stateful simulation.}
The working vector $\mathbf{w}$ is updated after each transfer. This
means suggestions are consistent: if CA~A transfers to CA~B, the
subsequent evaluation of CA~D will see B's updated load and will not
try to overload B again.

% ────────────────────────────────────────────────────────────
\section{Stage 3C --- Decision Engine}
\label{sec:decisions}
% ────────────────────────────────────────────────────────────

\subsection{Overview}

The decision engine translates numerical workload fractions into
prioritised, plain-language action recommendations. Each recommendation
has a priority level and a target (either the whole team or a specific CA).

\subsection{Priority Classification}

\begin{tcolorbox}[defbox, title={Priority Rules}]
Let $\bar{w}_0 = \frac{1}{C}\sum_c \mathbf{W}_{0,c}$ be the team average
and let $\hat{w}_c = \max_i \mathbf{W}_{i,c}$ be the peak workload
for CA $c$.

\medskip
\textbf{Team-level:}
\[
  \text{priority} = \begin{cases}
    \textcolor{highred}{\textbf{HIGH}}   & \text{if } \bar{w}_0 > 1.0 \\
    \textcolor{medorange}{\textbf{MEDIUM}} & \text{if } 0.85 < \bar{w}_0 \leq 1.0 \\
    \text{(no team alert)} & \text{otherwise}
  \end{cases}
\]

\medskip
\textbf{Per-CA:}
\[
  \text{priority}(c) = \begin{cases}
    \textcolor{highred}{\textbf{HIGH}}     & \text{if } \mathbf{W}_{0,c} > 1.0 \\
    \textcolor{medorange}{\textbf{MEDIUM}} & \text{if } \mathbf{W}_{0,c} \leq 1.0
                                             \text{ but } \hat{w}_c > 1.0 \\
    \textcolor{lowgreen}{\textbf{LOW}}     & \text{otherwise}
  \end{cases}
\]
\end{tcolorbox}

\textbf{HIGH (current overload)} signals that action is required now,
before the week is over.

\textbf{MEDIUM (forecast overload)} signals that the current week
is manageable, but a peak is predicted within the 78-week horizon.
The recommendation includes the exact date of the predicted peak.

\textbf{LOW} confirms that the CA is within safe capacity bounds
throughout the entire forecast horizon.

\subsection{Sorting}

The output list is sorted by priority order:
$\text{HIGH} \prec \text{MEDIUM} \prec \text{LOW}$,
ensuring the most urgent items appear first in the UI without any
client-side filtering.

% ============================================================
\chapter{API Layer}
% ============================================================

\section{Architecture}

The HTTP API is built with \textbf{FastAPI}, using asynchronous request
handling and automatic JSON serialisation via \textbf{Pydantic v2}
response models. All routes are mounted under the \texttt{/api} prefix.
Cross-Origin Resource Sharing (CORS) is configured to accept requests
from the Vite development server at \texttt{localhost:5173}.

\section{Server-Side State and Lazy Caching}
\label{sec:state}

The application is a single-user local tool. Rather than storing the
dataset in a database or on disk, the engine keeps a \textbf{module-level
singleton} in \texttt{backend/api/state.py}.

After a successful upload, the singleton holds:
\begin{itemize}
  \item \texttt{\_bundle} --- the \texttt{DataBundle} (parsed workbook data).
  \item \texttt{\_phase\_instances} --- the result of
        \texttt{build\_phase\_instances}, computed once and cached.
  \item \texttt{\_weekly\_forecast} --- the result of
        \texttt{forecast\_weekly}, computed once and cached.
  \item \texttt{\_monthly\_forecast} --- the result of
        \texttt{forecast\_monthly}, computed once and cached.
\end{itemize}

All four caches are reset to \texttt{None} on every new upload, forcing
recomputation from fresh data. Each derived cache is computed lazily:
only on the first API call that requires it.

\begin{tcolorbox}[notebox, title={Why lazy caching?}]
  Phase instances and the weekly forecast are the two most expensive
  computations (${\sim}300\,\text{ms}$ combined for 589 topics). By caching
  both after the first request, all subsequent endpoint calls are
  served in milliseconds from pre-computed in-memory DataFrames.
  Upload is the only operation that incurs the full computation cost.
\end{tcolorbox}

\section{Endpoint Reference}

\begin{table}[H]
\centering
\caption{API endpoints}
\begin{tabularx}{\linewidth}{@{} l l X @{}}
\toprule
\textbf{Method} & \textbf{Path} & \textbf{Description} \\
\midrule
\texttt{GET}  & \texttt{/api/health}    &
  Liveness check. Returns \texttt{\{"status":"ok"\}}. \\[4pt]
\texttt{POST} & \texttt{/api/upload}    &
  Accepts an \texttt{.xlsx} file (multipart/form-data), parses it,
  stores the bundle in the singleton, and returns summary metadata. \\[4pt]
\texttt{GET}  & \texttt{/api/overview}  &
  Returns team-level KPIs and a heatmap of workload fractions for
  the next $N$ weeks (\texttt{?weeks=26} default, max 78). \\[4pt]
\texttt{GET}  & \texttt{/api/team}      &
  Returns one row per active CA with current workload, peak workload,
  peak week, and active topic count. \\[4pt]
\texttt{GET}  & \texttt{/api/ca/\{id\}} &
  Returns full detail for one CA: individual forecast series, list
  of active phases with their VH, and all assigned topics. \\[4pt]
\texttt{GET}  & \texttt{/api/forecast}  &
  Returns the full forecast matrix. Supports \texttt{?view=weekly}
  (default) and \texttt{?view=monthly}. \\[4pt]
\texttt{GET}  & \texttt{/api/smoothing} &
  Runs the smoothing algorithm and returns the list of transfer
  suggestions. \\[4pt]
\texttt{GET}  & \texttt{/api/decisions} &
  Runs the decision engine and returns prioritised recommendations. \\
\bottomrule
\end{tabularx}
\end{table}

\section{Error Handling}

\begin{itemize}
  \item \texttt{HTTP 400} --- the uploaded file is not a \texttt{.xlsx} file.
  \item \texttt{HTTP 422} --- the workbook passes the format check but
        fails structural validation (missing sheet, missing column, etc.).
        The detail field contains a human-readable description.
  \item \texttt{HTTP 424} --- any analytical endpoint is called before a
        file has been uploaded. The response body explicitly instructs
        the client to upload first.
\end{itemize}

\section{Response Schemas}

All responses are typed with Pydantic v2 models, which provide:
\begin{itemize}
  \item Automatic JSON serialisation with field-level type coercion.
  \item Auto-generated OpenAPI documentation at \texttt{/docs}.
  \item Runtime validation that engine output matches the declared schema
        before it reaches the client.
\end{itemize}

% ============================================================
\chapter{Data Flow Summary}
% ============================================================

Figure~\ref{fig:dataflow} traces the complete data flow from user action
to rendered result.

\begin{figure}[H]
\centering
\begin{tcolorbox}[colback=gray!7, colframe=gray!50, sharp corners]
\small\ttfamily
\begin{verbatim}
  User uploads wtb_data.xlsx
        |
        v
  parse_workbook(bytes) --> DataBundle
        |
        v
  build_phase_instances(bundle)
    for each active CA x topic x phase:
      - compute VH = duration_hrs / delai_weeks
      - compute effective window [t_s*, t_e*]
      - classify is_stale, is_active_now
        |
        v
  forecast_weekly(bundle, phase_instances)
    for each week i in [w_0 ... w_H]:
      W[i, c] = sum(VH_p for p overlapping week i) / capacity
        |
        +--------> forecast_monthly  = resample("MS").mean()
        |
        +--------> n_active_topics   = count distinct active topics per CA
        |
        +--------> peak_analysis     = max over time axis
        |
        +--------> smoothing_suggestions (greedy transfer algorithm)
        |
        +--------> decisions (threshold-based priority classifier)
        |
        v
  FastAPI serialises via Pydantic --> JSON response to React frontend
\end{verbatim}
\end{tcolorbox}
\caption{End-to-end data flow}
\label{fig:dataflow}
\end{figure}

% ============================================================
\chapter{Testing}
% ============================================================

\section{Test Structure}

The backend test suite uses \textbf{pytest} and is organised in two
files under \texttt{backend/tests/}:

\begin{itemize}
  \item \texttt{test\_loader.py} --- validates the parsing layer:
        sheet detection, column presence, date coercion, config extraction.
  \item \texttt{test\_engine.py} --- validates the engine layer with
        25 unit tests across all major functions.
\end{itemize}

\section{Fixture Strategy}

All engine tests share a single \texttt{bundle} pytest fixture defined
in \texttt{conftest.py}, which loads the real
\texttt{docs/wtb\_data.xlsx} workbook. This means:
\begin{enumerate}
  \item Tests run against the actual production data schema, not a
        synthetic mock.
  \item Any structural change to the workbook is immediately caught.
  \item No test-specific data setup is required; the fixture is
        constructed once per test session.
\end{enumerate}

\section{Key Invariants Verified}

\begin{table}[H]
\centering
\caption{Selected test invariants}
\begin{tabularx}{\linewidth}{@{} l X @{}}
\toprule
\textbf{Function} & \textbf{Invariant tested} \\
\midrule
\texttt{build\_phase\_instances} &
  No inactive CA appears in output. \\
  &
  Every row has at least one non-null scheduled date. \\
  &
  Active-now rows never have excluded status. \\
  &
  Stale rows are never also active-now. \\
  &
  $t_s^* \geq t_0$ for all rows. \\
  &
  $t_e^* \geq t_s^*$ for all rows. \\
  &
  $\vh_p = \texttt{duration\_hrs}_p / \texttt{delai\_weeks}_p$ exactly. \\[4pt]
\texttt{forecast\_weekly} &
  Shape is $(H + 1) \times C$. \\
  &
  Columns match exactly the set of active CA identifiers. \\
  &
  All values are $\geq 0$. \\
  &
  All week index timestamps are Mondays (\texttt{weekday()} = 0). \\[4pt]
\texttt{forecast\_monthly} &
  All index timestamps have \texttt{day} = 1 (month-start). \\[4pt]
\texttt{decisions} &
  Priority values are in $\{\texttt{HIGH}, \texttt{MEDIUM}, \texttt{LOW}\}$. \\
  &
  Output is sorted HIGH $\prec$ MEDIUM $\prec$ LOW. \\
\bottomrule
\end{tabularx}
\end{table}

All 25 tests pass on the production dataset, confirming that every
algorithm invariant holds on real data.

% ============================================================
\chapter{Operational Notes}
% ============================================================

\section{Reproducibility}

The engine is fully reproducible by design. The reference date
\texttt{today} is read from the workbook's \texttt{config} sheet,
not from the system clock. Given the same \texttt{.xlsx} file,
any run on any machine at any time produces identical results.
This is essential for audit trails and management reporting.

\section{Capacity Interpretation}

A workload fraction of 1.0 represents one CA working at full capacity
(40 h/week). Values above 1.0 are analytically valid and expected for
overloaded CAs. The system does not cap these values; it reports them
as-is and flags them as requiring action.

\section{Catchup Window Sensitivity}

The catchup window $\Delta_c$ has the strongest effect on current-week
workloads. At $\Delta_c = 4$ weeks, an overdue phase contributes
$\vh_p \times 1$ to the current week's load. Increasing $\Delta_c$
spreads overdue work over more weeks and reduces current peaks;
decreasing it concentrates overdue work and increases peaks. This
parameter should be tuned to reflect the team's real operating cadence.

\section{Extensibility}

Because the engine is a set of pure functions with a well-defined input
(\texttt{DataBundle}) and output (DataFrames and plain dicts), it can
be extended without touching the API layer:

\begin{itemize}
  \item New phases: add rows to \texttt{phase\_params} and columns to
        \texttt{topics}; the engine loops over \texttt{PHASES} dynamically.
  \item New KPIs: add a new function to \texttt{engine.py} and call it
        from the relevant route.
  \item Alternative smoothing strategies (e.g.\ optimisation-based):
        replace or extend \texttt{smoothing\_suggestions} without
        modifying any other component.
\end{itemize}

% ============================================================
\chapter*{Glossary}
\addcontentsline{toc}{chapter}{Glossary}
% ============================================================

\begin{description}[leftmargin=3.5cm, labelwidth=3cm]
  \item[CA] Certification Agent --- a person responsible for qualifying
        a set of topics.
  \item[Topic] A supplier component or product that must go through
        the five-phase qualification process.
  \item[Phase] One of five sequential qualification steps:
        \textit{advanced, digital, static, dynamic, delivery}.
  \item[$\vh$] Weekly workload intensity of a phase, in hours per week.
        Defined as $\texttt{duration\_hrs} / \texttt{delai\_weeks}$.
  \item[Effective window] The $[t_s^*, t_e^*]$ interval used for workload
        attribution, accounting for overdue phases via the catchup window.
  \item[Catchup window] The $\Delta_c$-week period into which overdue
        work is compressed for forecast purposes.
  \item[Stale phase] A phase that ended more than $\Delta_s$ weeks ago
        without being formally closed; excluded from all forecasts.
  \item[Workload fraction] $\mathbf{W}_{i,c}$, the ratio of weekly hours
        demanded from CA $c$ to their capacity. 1.0 = full capacity.
  \item[DataBundle] The Python dataclass holding the four parsed DataFrames
        that flow through the engine.
  \item[VAL / CANCELED] Terminal phase statuses; work is definitively
        closed and excluded from all computations.
\end{description}

\end{document}
